\begin{textsample}{POS Dim 8 – human – Score 2.00 – mg/mg\_ef\_linguagens\_lp\_5\_p8.txt}  \label{ex:f8_pos_011}
O \textbf{Eixo} da Oralidade compreende as práticas de linguagem que ocorrem em situação oral com ou sem contato direto, face a face, como aula dialogada, webconferência, mensagem gravada, spot de campanha, jingle, seminário, debate, programa de rádio, entrevista, declamação de poemas ( com ou sem efeitos sonoros ), peça teatral, apresentação de cantigas e canções, playlist comentada de músicas, vlog de game, contação de histórias, diferentes tipos de podcasts e vídeos, \textbf{dentre} outras.
Envolve também a oralização de textos em situações socialmente significativas e interações e discussões envolvendo temáticas e outras dimensões linguísticas do trabalho nos diferentes campos de atuação.

% matched lemmas: dentrar, eixo
\end{textsample}
